\section{Introduction}
\label{sec:intro}

The Voting Rights Act of 1965 (\S 2) prohibits electoral practices that
result in minority voters having less opportunity than other members to
participate in the political process and elect representatives of their
choice.
In the redistricting context, Section 2 creates an affirmative obligation:
where a minority community satisfies the three-part \emph{Gingles} test
\citep{gingles1986}, mapmakers must provide the community with a realistic
opportunity to elect candidates of its choice.

Algorithmic redistricting systems face a structural tension here.
A map drawn by minimum-edge-cut bisection makes no reference to racial
or political data: it sees only boundary lengths and census populations.
This neutrality is legally valuable --- it defeats racial-predominance challenges
under \emph{Shaw v.\ Reno} \citep{shaw1993} --- but it means the algorithm
has no mechanism to detect or correct VRA shortfalls.

\paragraph{Why explicit minority constraints fail.}
The natural engineering response is to add minority population as a second
METIS constraint, alongside population balance.
B.3 \citep{b3paper} tested this approach (``MetisVra'') on five
majority-minority states with historically significant VRA litigation.
The results were damaging: MetisVra failed completely in 3 of 5 states
(Alabama, Louisiana, South Carolina), achieving zero majority-minority
(MM) districts across all 140 tested parameter configurations.

The failure mode is \textbf{constraint scale mismatch}.
METIS's multi-constraint solver treats all constraint dimensions with
the same numerical tolerance.
Population balance operates on vertex weights of 10,000--200,000 people,
requiring $\pm 0.5\,\%$ precision ($\pm 50$--$\pm 1000$ people).
Minority VAP fraction operates on vertex weights of 0.0--1.0, where
a $\pm 0.5\,\%$ tolerance is $\pm 0.005$ --- effectively zero.
The scale difference is 60--800$\times$.
METIS's local search cannot simultaneously satisfy both constraints:
relaxing the minority constraint to a useful tolerance (e.g., $\pm 50\,\%$)
provides so little guidance that METIS ignores it; tightening it creates
a numerical conflict that produces infeasible partitions.

\paragraph{VRASection: geographic alignment as tie-breaker.}
We propose a different approach grounded in Gingles Prong 1 itself.
Gingles requires that the minority community be ``sufficiently large
\emph{and geographically compact}'' to form a majority in a single-member
district \citep{gingles1986}.
Geographic compactness is not an obstacle to algorithmic redistricting:
it is a property that is \emph{already encoded in the census-tract adjacency
graph}.
A minority community that satisfies Prong 1 is, by definition, one whose
tracts cluster together in the adjacency graph.

VRASection makes this clustering measurable without adding a METIS constraint.
At the first bisection level, it evaluates each candidate ratio $i:(k-i)$
by a two-component score:

\begin{enumerate}
\item \textbf{Compactness}: the isoperimetrically normalised edge-cut
  $\mathrm{EC}_\text{norm}(i)$ from GeoSection \citep{b8paper} (lower is better).
\item \textbf{Geographic alignment}: $A(\text{split}) = |\mathrm{MVAP}_\text{frac}(L)
  - \mathrm{MVAP}_\text{frac}(R)|$, where $\mathrm{MVAP}_\text{frac}$ is the
  minority VAP fraction on each side of the bisection.
\end{enumerate}

The combined selection score is:
\[
  \text{score}(\text{ratio}) = \mathrm{EC}_{\text{norm}}
  - w_\text{vra} \cdot A \cdot \max(\mathrm{EC}_{\text{norm}},\, 1)
\]

Lower scores are preferred.
The alignment term \emph{subtracts} from the compactness score: a ratio
that places compact minority communities on one side of the cut receives
a lower (better) score than an otherwise identical ratio with dispersed
minority population.

\paragraph{Alabama: the lead case study.}
Applied to Alabama (7 districts, 27.5\,\% Black VAP, $w_\text{vra} = 0.40$),
VRASection selects a $2:5$ first-level bisection, separating the Black Belt
region (south-central Alabama, historically concentrated Black population)
from the northern industrial corridor and Gulf Coast.
GeoSection selects a $1:6$ split (Birmingham isolation --- the Alabama
caterpillar pathology).
The $2:5$ split more effectively satisfies Gingles Prong 1 because it
concentrates a larger fraction of Alabama's Black population on one side,
providing a geographic foundation for downstream MM district formation
without explicitly targeting race.

\paragraph{The legal distinction.}
VRASection uses minority VAP as input, which might superficially suggest
racial targeting.
The critical difference is in the function computed.
MetisVra's second constraint forces METIS to concentrate minority
population above a user-specified threshold --- race is a co-equal
optimization target.
VRASection's alignment score measures the \emph{spatial distribution}
of any population that the user designates as ``minority VAP''.
The algorithm cannot know whether the clustered population is Black,
Hispanic, Native American, or any other group.
It selects the bisection ratio that best exploits pre-existing geographic
concentration, which may or may not be a minority community.
At $w_\text{vra} = 0.40$, compactness is weighted $1 / (1 - 0.40) = 1.67\times$
more heavily than alignment, ensuring geographic compactness remains the
predominant factor under \emph{Shaw}'s racial-predominance test.

\paragraph{Contributions.}
This paper makes four contributions:

\begin{enumerate}
\item \textbf{VRASection algorithm} (\S\ref{sec:algorithm}): formal
  definition of the alignment score, modified objective function, and
  algorithm pseudocode, including legal analysis of why the design
  avoids racial predominance.

\item \textbf{Alabama case study} (\S\ref{sec:evaluation}): the $2:5$
  vs.\ $1:6$ comparison as the lead empirical result, with analysis of
  the Black Belt geographic concentration that drives the split.

\item \textbf{Comparison framework} (\S\ref{sec:evaluation}): a six-state
  comparison table covering the historically significant VRA states,
  with cells marked for runs pending full pipeline comparison.

\item \textbf{Legal and sensitivity analysis} (\S\ref{sec:discussion}):
  the Shaw predominance line as a function of $w_\text{vra}$, relationship
  to MetisVra failures, and comparison to the Callais evidence standard.
\end{enumerate}

\paragraph{Paper organisation.}
Section~\ref{sec:related} reviews related work on VRA compliance,
gerrymandering jurisprudence, and geographic concentration.
Section~\ref{sec:algorithm} defines VRASection formally.
Section~\ref{sec:evaluation} presents the Alabama empirical result
and the multi-state comparison framework.
Section~\ref{sec:discussion} discusses legal implications, $w_\text{vra}$
sensitivity, and limitations.
Section~\ref{sec:conclusion} concludes.
